\section{Development Prerequisites and Toolchain Setup}
\label{sec:development-prerequisites}

% Section 01.05: Development Prerequisites and Toolchain Setup

Before we can build and flash embedded Rust code to the EFR32MG24, we need to install and configure several tools. This section provides a complete setup guide for Linux, macOS, and Windows.

\subsection{Required Tools}
\label{sec:required-tools}

\textbf{Rust Toolchain}:

\begin{itemize}
    \item \textbf{Rust 1.70+} with support for the thumbv8m.main-none-eabihf target
    \item \textbf{cargo-binutils} for examining compiled binaries
    \item \textbf{cargo-size} for checking binary size
    \item \textbf{cargo-objdump} for disassembly
\end{itemize}
\subsection{Installation on Linux (Ubuntu/Debian)}
\label{subsec:install-linux}
\begin{minted}{bash}
# Install Rust via rustup
curl --proto '=https' --tlsv1.2 -sSf https://sh.rustup.rs | sh

source $HOME/.cargo/env

# Add ARM Cortex-M33 target
rustup target add thumbv8m.main-none-eabihf

# Install cargo tools
cargo install cargo-binutils
rustup component add llvm-tools-preview

# Install system dependencies
sudo apt-get update
sudo apt-get install -y \
    build-essential \
    pkg-config \
    libssl-dev \
    git \
    wget
\end{minted}
\subsection{Installation on macOS}
\label{subsec:install-macos}
\begin{minted}{bash}
# Install Rust via rustup
curl --proto '=https' --tlsv1.2 -sSf https://sh.rustup.rs | sh

source $HOME/.cargo/env

# Add ARM Cortex-M33 target
rustup target add thumbv8m.main-none-eabihf

# Install cargo tools
cargo install cargo-binutils
rustup component add llvm-tools-preview

# Install Homebrew if not already installed
/bin/bash -c "$(curl -fsSL https://raw.githubusercontent.com/Homebrew/install/HEAD/install.sh)"

# Install system dependencies
brew install git wget
\end{minted}
\subsection{Installation on Windows}
\label{subsec:install-windows}
\begin{enumerate}
    \item Download and install from https://rustup.rs/
    \item Run the installer with default settings
    \item After installation, open a new PowerShell/Command Prompt
    \item Install the ARM target:
\end{enumerate}

\begin{minted}{bash}
rustup target add thumbv8m.main-none-eabihf
cargo install cargo-binutils
rustup component add llvm-tools-preview
\end{minted}

\subsection{Hardware Debug Probe Setup}
\label{sec:debug-probe-setup}

The EFR32MG24 requires a debug probe for flashing and debugging. We'll use probe-rs, a Rust-based tool.

\subsection{Installing probe-rs}
\label{subsec:install-probe-rs}
\begin{minted}{bash}
# Install probe-rs
cargo install probe-rs --features cli

# Verify installation
probe-rs list
\end{minted}

Expected output (when hardware is connected):

\begin{minted}{text}
Probing OpenUSB... Done
[0]: J-Link Silicon Labs (<serial-number>) (J-Link)
\end{minted}

\subsection{Flashing: one sanctioned path}
\label{subsec:install-silabs-tools}

Silicon Labs' Simplicity Commander exists, and it is the right tool for one job: the
authorisation that unlocks a secured part. It is \emph{not} the flashing path this book
uses, because a bare \codeinline{commander flash} will happily write wherever the image
says — including over the bootloader.

What this project uses instead, and why it is worth the extra script:

\begin{minted}{bash}
# 1. show where the image would land, without touching the board
python3 scripts/elf-load-range.py target/thumbv8m.main-none-eabihf/release/examples/pac_imu
# load_min=0x08006000 load_max=0x08006813 size=2068 segments=3

# 2. dry run: prints the plan and refuses anything with load_min < 0x08006000
bash scripts/flash-lnx.sh "$ELF"

# 3. write, then verify by reading the flash back
bash scripts/flash-lnx.sh "$ELF" --confirm
\end{minted}

The board hangs off a separate host (\codeinline{lnx}) reachable over SSH, which is why
the script copies the ELF over and drives the programmer there. If your board is on the
machine you build on, the same script works locally — the guards, not the transport, are
the point.

\begin{quote}
\textbf{Measured caveat.} \codeinline{probe-rs} works for the first attach and then
desynchronises on this board's onboard CMSIS-DAP bridge
(\codeinline{Command ID in response (0x14) does not match sent command ID}); a USB
reset does not clear it and only a physical replug does. OpenOCD talks to the same probe
without the problem, so the script uses OpenOCD for reads and \codeinline{probe-rs} for
the write-and-verify step.
\end{quote}

\subsection{Updating Debug Probe Firmware}
\label{subsec:update-firmware}

If using a J-Link probe:

\begin{minted}{bash}
# On Linux/macOS
JLinkExe -CommanderScript <(echo "exit")

# Follow prompts to update firmware
\end{minted}

\subsection{Development Environment Setup}
\label{sec:dev-env-setup}

\subsection{VS Code Configuration}
\label{subsec:vscode-setup}

Install these extensions:
\begin{itemize}
    \item rust-analyzer (official Rust support)
    \item Even Better TOML (for Cargo.toml files)
    \item CodeLLDB (for debugging)
\end{itemize}

Create \filepath{.vscode/settings.json}:
\begin{minted}{json}
{
    "rust-analyzer.cargo.target": "thumbv8m.main-none-eabihf",
    "rust-analyzer.checkOnSave.allTargets": false,
    "rust-analyzer.cargo.features": ["rt"],
    "rust-analyzer.procMacro.ignored": {
        "cortex-m-rt": ["interrupt"]
    }
}
\end{minted}

\subsection{Cargo Configuration for Embedded}
\label{subsec:cargo-config}

Create \filepath{.cargo/config.toml}:
\begin{minted}{toml}
[target.thumbv8m.main-none-eabihf]
rustflags = [
    "-C", "link-arg=-Tlink.x",
    "-C", "link-arg=-Tdefmt.x",  # If using defmt
]

runner = "probe-rs run --chip EFR32MG24B220F1536IM48"

[build]
target = "thumbv8m.main-none-eabihf"

[env]
DEFMT_LOG = "debug"
\end{minted}

\subsection{Git Configuration}
\label{subsec:git-config}

Set up Git for embedded development:
\begin{minted}{bash}
git config --global user.name "Your Name"
git config --global user.email "your.email@example.com"

# Set up LFS for large files (like SVD files)
git lfs install
\end{minted}

Create \filepath{.gitattributes}:
\begin{minted}{text}
# Track large files with Git LFS
*.svd filter=lfs diff=lfs merge=lfs -text
*.pack filter=lfs diff=lfs merge=lfs -text
*.a filter=lfs diff=lfs merge=lfs -text
\end{minted}

\subsection{Verifying Your Setup}
\label{sec:verify-setup}

Run this verification script to ensure everything is correctly installed:
\begin{minted}{bash}
#!/bin/bash

echo "=== Checking Rust Installation ==="
rustc --version
cargo --version

echo ""
echo "=== Checking Targets ==="
rustup target list | grep "thumbv8m.main-none-eabihf (installed)"

echo ""
echo "=== Checking Cargo Tools ==="
cargo size --version 2>/dev/null && echo "✓ cargo-size installed"

echo ""
echo "=== Checking probe-rs ==="
probe-rs --version

echo ""
echo "=== Checking debug probe ==="
probe-rs list
\end{minted}

Expected output:
\begin{minted}{text}
=== Checking Rust Installation ===
rustc 1.70.0 (90c541806 2023-05-31)
cargo 1.70.0 (90c541806 2023-05-31)

=== Checking Targets ===
thumbv8m.main-none-eabihf (installed)

=== Checking Cargo Tools ===
✓ cargo-size installed

=== Checking probe-rs ===
probe-rs 0.18.0

=== Checking debug probe ===
[0]: J-Link Silicon Labs (xxxx) (J-Link)
\end{minted}

If all checks pass, you're ready to start development!

\begin{tipbox}
    \textbf{Common Issues}:
    \\• probe-rs list shows nothing: Check USB cable, ensure probe has power
    \\• Target not installed: Run rustup target add thumbv8m.main-none-eabihf
    \\• cargo-size not found: Run cargo install cargo-binutils
    \\• Permission denied (Linux): Add your user to dialout or plugdev group
\end{tipbox}

\begin{exercise}[Easy]
    Create a shell script that automates the entire setup verification process for your team.
\end{exercise}

\begin{exercise}[Medium]
    Write a cargo subcommand that checks for required tools and provides helpful installation instructions if anything is missing.
\end{exercise}

\begin{exercise}[Hard]
    Using Docker, create a container with the complete development environment pre-installed so new team members can start immediately.
\end{exercise}
\clearpage
